\documentclass[11pt]{article}
\usepackage[T1]{fontenc}
\usepackage{textcomp}
\usepackage[margin=0.75in]{geometry}
\usepackage{amsmath,amssymb,amsthm}
\usepackage{array,booktabs}
\usepackage{hyperref}
\setlength{\parindent}{0pt}
\newtheorem{theorem}{Theorem}
\theoremstyle{definition}
\newtheorem{definition}{Definition}
\title{Flow Proof Artifact: Finset-card-inter-union-diff-empty}
\author{Generated by \texttt{flow doc proof}}
\date{Flow Proof Book}
\begin{document}
\maketitle
Intersection union difference empty has bounded cardinality.\\[0.5em]
\textbf{Source.} graham-knuth-patashnik — *Concrete Mathematics*\\[0.5em]
\bigskip
\noindent{\large \textbf{Derived fact 1.} \textit{card((s ∩ t) ∪ (s \textbackslash\{\} t)) <= card(s ∪ t)} \hfill Finset \cdot cardinality \cdot inter union diff empty card bound}\par
\smallskip\noindent\textit{Coordinate: Finset \cdot cardinality \cdot inter union diff empty card bound}\par
\smallskip\noindent\textit{Source: graham-knuth-patashnik}\par
\smallskip\noindent\textit{Built on: inter union diff empty card derived, for cardinality on Finset, union diff inter empty card bound, for cardinality on Finset}\par
\medskip
\noindent\textbf{Goal.} card((s ∩ t) ∪ (s \textbackslash\{\} t)) <= card(s ∪ t)\par
\noindent$\displaystyle \forall s \in Finset \forall t \in Finset\quad card((s ∩ t) ∪ (s \ t)) \le card(s ∪ t)$\par
\medskip
\noindent
\begin{tabular}{@{} >{\raggedright\arraybackslash}p{0.06\textwidth} >{\raggedright\arraybackslash}p{0.47\textwidth} >{\raggedleft\arraybackslash}p{0.41\textwidth} @{}}
\textbf{\#} & \textbf{Proof} & \textbf{Mathematics} \\
\midrule
\textbf{1.} & We prove that inter union diff empty card bound for cardinality on Finset. & \\[0.45em]
\textbf{2.} & We invoke the derived fact governing cardinality on Finset: inter union diff empty card derived, for cardinality on Finset (instantiated for s, t). & \\[0.45em]
\textbf{3.} & We invoke the derived fact governing cardinality on Finset: union diff inter empty card bound, for cardinality on Finset (instantiated for s, t). & \\[0.45em]
\textbf{4.} & From step 2 and step 3, this implies card((s ∩ t) ∪ (s \textbackslash\{\} t)) is at most card(s ∪ t). Hence proven. & $\displaystyle card((s ∩ t) ∪ (s \ t)) \le card(s ∪ t)$ \\[0.45em]
\bottomrule
\end{tabular}
\label{thm:Finset-cardinality-inter_union_diff_empty_card_bound}
\par\medskip\hrule
\end{document}
